\Introduction

Современное состояние мировой микроэлектроники характеризуется интенсивным развитием архитектур микроконтроллеров и растущими требованиями к разработки встроенных систем. В условиях геополитической нестабильности и технологических ограничений особую актуальность приобретает создание отечественных решений в области микроэлектроники, способных обеспечить технологическую независимость критически важных отраслей экономики.

Микроконтроллер MIK32 демонстрирует потенциал отечественных технологий в создании конкурентоспособных решений. Однако использование потенциала данного микроконтроллера требует создания специализированного программного обеспечения, которое обеспечит удобство разработки и оптимальное использование аппаратных ресурсов.

Актуальность исследования обусловлена несколькими факторами. Во-первых, растущая потребность в импортозамещении микроэлектронных компонентов требует не только создания отечественной элементной базы, но и разработки соответствующего программного инструментария. Во-вторых, открытая архитектура RISC-V предоставляет уникальные возможности для создания оптимизированных решений без ограничений проприетарных лицензий. В-третьих, отсутствие зрелых инструментов разработки для MIK32 создает барьеры для широкого внедрения данного микроконтроллера в промышленных применениях.

Анализ современного состояния разработки встроенных систем показывает, что фреймворки играют важную роль в повышении производительности разработчиков и качества создаваемых решений. Существующие фреймворки и операционные системы реального времени для микроконтроллеров, такие как STM32CubeHAL для семейства STM32, Mbed OS, FreeRTOS, Zephyr или ESP-IDF для ESP32, демонстрируют преимущества в скорости разработки и надежности результирующих систем. Однако специфика архитектуры MIK32 и требования отечественного рынка обусловливают необходимость создания специализированного решения.

Объектом исследования является микроконтроллер MIK32 (К1948ВК018) как представитель класса 32-битных микроконтроллеров с архитектурой RISC-V, включая его аппаратную архитектуру и периферийные устройства.

Предметом исследования выступают методы и принципы проектирования программных фреймворков для микроконтроллеров, включая архитектурные паттерны абстракции аппаратных ресурсов, подходы к созданию универсальных программных интерфейсов и механизмы обеспечения переносимости программного кода.

Цель работы заключается в упрощении процесса разработки встроенных приложений за счет исследовании и создания прототипа фреймворка для микроконтроллера MIK32, который обеспечит абстракцию аппаратных ресурсов.

Для достижения поставленной цели определены следующие задачи:

\begin{enumerate}
	\item Выполнить сравнительный анализ существующих фреймворков и операционных систем реального времени для микроконтроллеров с целью выявления лучших практик и архитектурных решений.
	\item Провести анализ особенностей микроконтроллера MIK32, включая изучение процессорного ядра SCR1, периферийных устройств и специфики программной модели RISC-V.
	\item Разработать прототип многоуровневого фреймворка, обеспечивающую разделение между аппаратно-зависимыми и аппаратно-независимыми компонентами.
	\item Спроектировать и реализовать унифицированные подсистемы управления периферийными устройствами.
	\item Провести апробацию фреймворка в реальных промышленных проектах различной сложности.
\end{enumerate}

Методы исследования включают: анализ архитектуры микроконтроллеров и существующих программных решений; экспериментальное моделирование и тестирование на реальных аппаратных платформах; сравнительный анализ с существующими решениями.

Научная новизна работы заключается в разработке модульной архитектуры фреймворка, которая учитывает требования российского рынка встренных стистем.
Практическая значимость исследования определяется успешной апробацией фреймворка в промышленных проектах различной сложности, подтвердившей его универсальность и практическую применимость. Формированием методологической базы для создания аналогичных фреймворков для других отечественных микроконтроллеров. Созданием готового к использованию инструментария для разработки встроенных систем на базе MIK32, что способствует ускорению процесса внедрения отечественного микроконтроллера в промышленные применения.

\Abbrev{API}{Application Programming Interface — интерфейс прикладного программирования}
\Abbrev{HAL}{Hardware Abstraction Layer — слой абстракции аппаратного обеспечения}
\Abbrev{GPIO}{General Purpose Input/Output — порты ввода-вывода общего назначения}
\Abbrev{I2C}{Inter-Integrated Circuit — двухпроводной последовательный интерфейс}
\Define{Фреймворк}{программная платформа, определяющая структуру программной системы и облегчающая разработку приложений}
